\subsection{Prepared charge localization and relative spin phase}

The complete charged energy measurement also controls position for its
specified preparation.  Write $H=H_0+\lambda T+4N_{\rm ref}$ for the
eight-cell interacting source, with $\lambda=2^{-34}$ and
$\|\lambda T\|\leq v=360\,2^{-34}$.  The reference is a retained static
charge at the origin.  Let $\chi=P_I(H)B_r\Psi$ be the unnormalized
prepared charged line, and let $\phi$ be its atomic comparison vector.
The complete spectral-projector and preparation bounds give
$\|\chi-\phi\|\leq\epsilon$, $H_0\phi+4N_{\rm ref}\phi=a\phi$,
and $\Pi_0\phi=\phi$.  In the physical one-reference sector, $\Pi_y$
is the single-charge-at-$y$ projector: it requires every other matter cell
to be a gauge singlet and retains all link states.  These projectors are
pairwise orthogonal on that sector.  If $w=\|\phi\|^2\in[w_-,w_+]$ and
$p_-=(\sqrt{w_-}-\epsilon)^2>0$, Duhamel's identity gives
\[
 \Pr(\Pi_0\mid\text{line},t)\geq
 1-\min\left\{1,
 \frac{(\epsilon+|t|v\sqrt{w_+})^2}{p_-}\right\}.
\]
The complete cloud effect is $\Pi_{\rm cloud}=I-\sum_y\Pi_y$, so
$I-\Pi_0=\Pi_{\rm cloud}+\sum_{y\ne0}\Pi_y$ retains every remote
single-charge outcome and every multi-charged-cell cloud outcome,
including screening configurations.  Neutral local excitations and link
excitations need not lie outside $\Pi_0$.  For the native $e^c$ removal channel the
origin probability exceeds $0.9999999127$ at preparation and
$0.9995481684$ at $t=10^6$ in source units; the latter interval exceeds
$110{,}036$ atomic source-gap cycles.  These are bounds on the complete
prepared state under its actual Hamiltonian.  They do not apply to every
state in the same energy bin, and the fixed compensator is part of the
construction.

The source's central spin rotation has an exact finite readout on these
same states.  Its native generator is
$S_z=(N_\uparrow-N_\downarrow)/2$, hence
$e^{-2\pi i S_z}=(-1)^N$ on the full source Fock space.  The prepared
neutral cube is even, each successful elementary removal is odd, and
the full Hamiltonian and energy filter preserve that parity.  Controlled
source rotation and a retained coherent pointer therefore distinguish
the minus sign at $2\pi$ from the identity at $4\pi$.  Rotating the odd
compensating reference as well cancels the sign and supplies a distinct
control.  All preparation, pump and energy-bin failures are retained.
This is a finite source and apparatus result; it does not identify a
continuum electron angular momentum or derive an apolar tensor field.

\subsection{Finite electromagnetic and matter-wave readouts}

For a finite incidence complex, let $D$ be the gradient and $C$ the curl, with $CD=0$. Positive electric and magnetic energy, Faraday evolution and the declared time reversal give
\[
 \dot E=c^2C^*B,\qquad \dot B=-CE.
\]
A total-occupation cutoff preserves the linear Heisenberg equations when the positive boundary completion is retained. The energy differences of $H=\sum_j\omega_j a_j^*a_j$ establish $[H,a_j]=-\omega_j a_j$ and hence the linear Heisenberg equations on every retained state, including the cutoff boundary. This does not assert canonical coordinate commutators on the finite carrier. Propagation therefore reduces to the coordinate matrices, without diagonalizing the full Fock carrier.

For the declared periodic cubic extension, the two transverse polarizations obey
\[
 \omega_a^2=\frac{4c^2}{a^2}\sum_j\sin^2\frac{ak_j}{2},\qquad
 0\le c^2|k|^2-\omega_a^2\le\frac{c^2a^2}{12}\sum_jk_j^4.
\]
Independent real-space evolution and exact modal propagation test this bound. The same gradient supplies the minimum-energy neutral electrostatic solution
\[
 E_L=D(D^*D)^+\rho,\qquad
 U=\frac{g^2}{2}\rho^*(D^*D)^+\rho.
\]
Potential differences and forces approach the Coulomb form under the stated volume and spacing controls. The coefficient $g^2/(4\pi)$ retains the declared field normalization. A Dirichlet consumer with supplied nonrelativistic kinetic energy has negative ground and excited levels; exact reflection sectors reduce the largest $95^3$-site source to $110{,}592$ even-sector rows. The ground-to-odd spacing decreases from approximately $0.686$ to $0.4204$ under the retained refinement, compared with the hydrogenic reference $0.375$ in its declared units. This is a conditional bound-state calculation with unresolved continuum and finite-volume errors, not a precision hydrogen spectrum.

A separate complete sixteen-row native mass-channel calculation retains all eight valleys and verifies the massive second-order wave law through direct time evolution. Its finite-spacing boost defect is bounded and decreases under refinement. The supplied mass coefficients, electromagnetic normalization and physical particle identification are separate inputs; these finite calculations do not identify a common interacting relativistic source.

\subsection{Metric response and conditional cosmological consequences}

For the reciprocal scalar action density
$C(r)=\operatorname{diag}(-1/r,1/r,r^2,r^2)$, the metric dictionary
$g=\sqrt{-\det C}\,C^{-1}$ gives unit proper spatial volume, two-face area
$A=2/r$, and lapse $L=r$.  On the two branches $r=9/10,11/10$, a complete
matter-randomization intervention followed by the unchanged thermal update
therefore obeys
\[
 \Delta\langle L\rangle=-\frac{99}{200}\Delta\langle A\rangle.
\]
The lapse response per unit randomization is enclosed about
$0.0003103930363$, retaining all sixteen joint levels and every rejected
proposal.  This is a finite metric readout; no spatial potential or
frequency comparison between separated clocks is inferred.

The conditional regional Einstein bound has a direct weak-field consumer.
For a static spherical pressureless source, with negligible or subtracted
cosmological background, let every component of the Einstein residual be
bounded by $\epsilon_E$ and vanish outside radius $R$.  In
$ds^2=-(1+2\Phi)dt^2+(1-2\Psi)d\mathbf x^2$, direct curvature
differentiation gives
$\nabla^2\Phi=4\pi G_N\rho+(E_{00}+\sum_iE_{ii})/2$.
Regularity at the origin and vanishing potential at infinity imply
\[
 |\delta\Phi'(r)|\le\frac{2\epsilon_E R^3}{3r^2},
 \qquad |\delta\Phi(r)|\le\frac{2\epsilon_E R^3}{3r},\qquad r\ge R.
\]
Thus the matched regional entropy and area errors bound the exterior
force directly.  The geometric and stress identifications required by
that regional theorem, together with the weak-field and compact-support
assumptions, are explicit premises of this implication.

For the separately specified capacity-density branch
$\rho_{\rm DE}=\kappa/N(a)$, Friedmann continuity yields
$w=-1+\tfrac13d\ln N/d\ln a$.  A flat comparison history with
$E=H/H_0$ satisfies
\[
 \frac{N(a)}{N(1)}=
 \frac{\Omega_{\rm DE}}
 {E(a)^2-\Omega_m a^{-3}-\Omega_r a^{-4}}.
\]
This normalization-free identity tests capacity evolution while leaving
its amplitude unselected.  A numerical cosmological constant requires
the additional global-capacity and horizon identifications; the displayed
identity does not determine them.

\subsection{Magnetic response of the native kinetic family}

The complete sixteen-row native kinetic matrices obey the Clifford relations. In a declared Peierls gauge extension with scalar mass $m$, their finite Hamiltonian satisfies
\[
 H^2=m^2+\Pi_x^2+\Pi_y^2+\Sigma_z i[\Pi_x,\Pi_y],\qquad
 \Sigma_z=-i\gamma_x\gamma_y,\quad S_z=\Sigma_z/2.
\]
A smooth uniform external field therefore gives the nonrelativistic spin term $qBS_z/m$, corresponding to tree-level $g=2$ without an inserted Pauli coefficient. Exact internal projectors retain all eight valley labels and reduce the largest $266{,}256$-row example to two $16{,}641$-row spatial operators. Tapered finite-packet energy responses give $g_{\rm eff}=1.95802,1.98942,1.99734$ at spacings $1,1/2,1/4$; coherent spin-phase measurements agree, with an explicit square-root truncation bound. The gauge extension, mass and external field are supplied. No self-sourced dipole, physical electron identification or radiative anomalous moment is inferred from this external-field calculation.

\subsection{Conserved charge without a winding dictionary}

The native internal current $Q=T_3+Y$ and colour Cartan generator $C$, induced by $\operatorname{diag}(1,1,-2)/3$, obey $3Q_i\in\mathbb Z$ and $Q_i+C_i\in\mathbb Z$ on all fifteen internal weights. Exterior products lift these identities to every fermion number and any number of spatial or family copies:
\[
 \operatorname{spec}(d\Gamma(Q))\subset\tfrac13\mathbb Z,
 \qquad e^{2\pi i d\Gamma(Q)}\Gamma(Z_{\rm colour})=I.
\]
Colour-singlet charge eigenvalues are integral, and a half-unit charge eigenvalue is absent. Every specified guarded texture-record map is $\sqrt{p_r}\Gamma(U_r\otimes I)$ and preserves these charge sectors. A mixture can have mean charge one-half while its retained outcomes are zero and one. This is finite spectral quantization and branchwise charge conservation, conditional on the declared current, normalization and record law; it does not exclude charge-coherent states or other external interventions.

The compact texture independently has integer shell degrees $-1,+1$. Every accepted guard-preserving update is a nonsingular homotopy of its closed shell; rejected updates are identity maps. Unit degree therefore cannot split into two equal integer degrees under that law. The same fixed texture nevertheless admits neutral and charged fermionic occupations, so the current cannot be a function of winding alone on the unchanged carrier. The direct current proof does not require that false identification. An observed electromagnetic dictionary and a local physical dressing require their own source attachment.

\subsection{Native scalar attraction and its universality obstruction}

In units of the local factor $\sqrt{14}$, the three-oscillator reciprocal source has
\[
 M(r)=a(r)D-\frac{C}{7r},\qquad
 a(r)=\frac r4+\frac{1}{7r}+\frac{3r^2}{28},\qquad
 J=\left.\frac{dM}{dr}\right|_{r=1}=\frac9{28}D+\frac17 C.
\]
Here $D_{ss}=2\operatorname{popcount}(s)+3$ and $C=XXI+IXX$. A specified matter-state contrast supplies charge $j=\operatorname{Tr}(\delta\rho J)$. On a declared spatial extension with positive elastic energy $\kappa q^TLq/2+j^Tq$, completing the square yields the mutual interaction $-j_xj_y(L^{-1})_{xy}/\kappa$. The local coupling comes from the source derivative; spatial gluing, stiffness, zero boundary and a massless geometry mode are separate architecture inputs.

A finite dyadic geometry carrier admits exact minimum-energy enclosures from a trial residual and a lower bound on the Dirichlet Laplacian. For the occupation contrast, the native nonlinear remainder is exactly
\[
 \frac{q^2(7+3q)}{14(1+q)}.
\]
Bounding it on the entire finite window controls the nonlinear minimum without solving a nonlinear optimization. The retained separation-six force interval is strictly attractive. However, two certified preparations above the actual local ground have positive excess energies and opposite $J$ contrasts. This scalar extension consequently does not supply universal positive-mass gravity or a derived Newton constant.
